\subsection{Standardised measures of interest rate risk in the banking book}
\label{sec:rw-standard}

The Basel Committee's standard for interest rate risk in the banking book defines two measures, the change in the economic value of equity under six prescribed rate scenarios and the change in net interest income under two parallel scenarios, each computed by currency in which the bank has material positions \citep{BCBS_d368,BCBS_SRP31}. It assumes that the rate shock of a currency is a function of that currency's own history, gives the supervisor's reading of the largest absolute rate moves, and adds the currencies by summing the positive parts of the losses, so that a gain in one currency never offsets a loss in another. The 2024 recalibration replaced the global factors of the 2016 table by currency-specific shocks, the 99.9th percentile of the absolute changes in the average rate over rolling six-month windows from 2000 to 2023, rounded to 25 basis points, and left the shapes, the buckets and the floor provision unchanged \citep{BCBS_d578}. Its table lists 21 currencies, and the standard lets jurisdictions recalculate the shocks of their domestic currencies and lets supervisors set a floor on post-shock rates that is not greater than zero. What the standard gives is a comparable measure across banks; where it fails for a Chilean bank is that neither the peso nor a unit of account indexed to prices appears in the table, so the standard cannot be applied without a national text.

The national text is Annex 1 of chapter 21-13 of the compilation of rules of the Comisi\'on para el Mercado Financiero \citep{CMF_RAN_21_13}. It keeps the six scenarios and the bucket grid, fixes the shock sizes of the peso and of the UR in a table of its own, floors the post-shock rates at zero for the peso and at minus two percent for the UR, treats the indexed flows as one currency, combines the two legs by the pair rule, and adds the indexation charge to the income measure. It also states the thresholds above which a bank is prioritised in the supervisory review. The supervisor's text is the specification of this paper, not a method it proposes. The review of the Federal Reserve of the supervision of Silicon Valley Bank is the closest public statement of why a standardised measure matters: it describes the bank's senior leadership as having failed to manage basic interest rate and liquidity risk \citep{FRB_SVB}. It says nothing about indexed units, and it is cited here only for the consequence of an understated exposure.

\subsection{Indexed units of account and real term structures}
\label{sec:rw-indexed}

Shiller treats an indexed unit of account as a money analogue that expresses prices, whose purchasing power is defined by an index and that is not used as a medium of exchange, and discusses the UF as the first successful example \citep{Shiller1998}. That work assumes the unit exists and asks why it is useful; it gives the monetary-theory reason that a flow in the unit is protected from inflation and does not give a risk measure for a bank that holds it. The pricing literature gives the instrument that is closest to a UF flow, the inflation-indexed government bond. Jarrow and Yildirim price such bonds and their derivatives in a Heath, Jarrow and Morton framework, fitting a three-factor model to the consumer price index and to the real and the nominal zero-coupon curves, in which the index converts a real value into a nominal one \citep{Jarrow2003,Heath1992}; a separation of the real and the nominal side of a flow comparable to the one that Annex 1 makes administrative when it names the UR. G\"urkaynak, Sack and Wright fit a smoothed yield curve to the indexed bonds of the United States Treasury and compare it with the nominal curve to obtain inflation compensation, and they show that the compensation is not a pure measure of expected inflation because it contains risk and liquidity premia \citep{Gurkaynak2010}. Ang, Bekaert and Wei estimate the components of the nominal curve with a regime-switching model and find a real curve that is fairly flat on average and an inflation risk premium that rises with maturity \citep{Ang2008}. Campbell, Shiller and Viceira describe the history of the indexed bond markets of the United States and the United Kingdom and the large fall in long real rates until 2008 \citep{Campbell2009}. Fleckenstein, Longstaff and Lustig show that a Treasury bond and a replicating inflation-swapped indexed bond can differ in price, so that the two curves are not tied by arbitrage in practice \citep{Fleckenstein2014}. These works fail, for the present purpose, in being about market prices of indexed securities and not about the rule that a supervisor applies to a bank's repricing flows; they supply the reason, used in Section~\ref{sec:method-curve}, that an indexed flow must be discounted on a real curve and the gap between the two curves is a breakeven rate, not a constant.

\subsection{Yield-curve models behind a zero curve}
\label{sec:rw-curves}

The zero curves of the experiments are fitted, not observed, and the choice of the curve model enters the measures. Nelson and Siegel propose a parsimonious three-factor form for the yield curve that fits level, slope and curvature \citep{Nelson1987}, and Svensson extends their functional form to make it more flexible, with Sweden over 1992 to 1994 as the example \citep{Svensson1994}. Diebold and Li make the three factors dynamic and forecast the term structure with them \citep{Diebold2006}. Litterman and Scheinkman study the common factors that affect bond returns \citep{Litterman1991}; their text was not available for this work and the characterisation rests on its title and on later descriptions of it as a factor analysis of a small number of common factors of the curve, which is the practice that a measure built on parallel, short and long shocks also follows. G\"urkaynak, Sack and Wright publish the yield curve of the United States Treasury from 1961 with a parametric form of this family \citep{Gurkaynak2007}. Models of the short rate and of the forward curve, those of Vasicek, of Hull and White and of Heath, Jarrow and Morton, generate the paths on which an earnings-at-risk measure is computed \citep{Vasicek1977,Hull1990,Heath1992}; the case behind this paper computes such measures too, but the standardised measures of this paper need only the zero curve at the 19 bucket midpoints. These models assume that the curve is a smooth function of maturity and give a curve that the regulation then shocks; they fail where the observed yields are few, and the peso curve of the case has three published tenors, so that the engine adds the policy rate as a fourth point (Section~\ref{sec:design-data}).

\subsection{Evidence on how banks bear interest rate risk}
\label{sec:rw-banks}

A second literature measures the exposure of real banks rather than defining a regulatory measure. Drehmann, Sorensen and Stringa build a framework that models the whole portfolio of a bank, with the repricing characteristics of its assets, liabilities and off-balance-sheet items, to measure the joint effect of credit and interest rate risk on economic value and capital adequacy, and apply it to a hypothetical but realistic bank under stress; it assumes a stylised balance sheet and shows that the two risks have to be measured together, not what a standardised rule adds to a measure \citep{Drehmann2010}. Memmel replicates the banking-book interest rate risk of banks with passive investment portfolios and finds that the changes of market value and of net interest income are highly correlated whatever the portfolio, while the portfolio composition has a large effect on the ratio of the two, which matters for the interpretation of stress tests; the two perspectives are the ones that the standard measures also use, and the study has no indexed unit \citep{Memmel2011}. Entrop, Memmel, Ruprecht and Wilkens extend the Ho and Saunders model of interest margins with interest rate risk and the expected return of maturity transformation, on German universal banks over 2000 to 2009, and find that all banks charge for macroeconomic interest rate volatility while the exposure from maturity transformation is priced only by the smaller savings and cooperative banks; it explains margins, not a regulatory measure \citep{Entrop2015}. Esposito, Nobili and Ropele measure the exposure of Italian banks during the crisis with the standardised duration gap of the Basel Committee and find that on-balance-sheet risk and derivatives substitute for each other, with a negative correlation between interest rate risk and liquidity risk; the duration gap is a first-order measure on one currency \citep{Esposito2015}. Hoffmann, Langfield, Pierobon and Vuillemey use novel data on the European banking sector and find that the exposure to interest rate risk is small on aggregate but heterogeneous, that net worth rises with rates for about half of the institutions, that loan-rate fixation conventions for mortgages drive the cross-section, and that banks hedge part of the exposure with derivatives \citep{Hoffmann2019}. Drechsler, Savov and Schnabl argue that maturity transformation hedges the interest rate risk of the deposit franchise instead of creating it \citep{Drechsler2021}. English, Van den Heuvel and Zakrajsek identify the reaction of bank stock prices to the interest rate movements that follow monetary policy announcements and find that prices fall after unanticipated increases in the level or the slope of the yield curve, more for banks that rely on core deposits and less for banks with a large maturity mismatch; the exposure is measured on prices and not by a supervisory rule \citep{English2018}, and Begenau, Piazzesi and Schneider measure interest rate and credit risk exposures in the fixed-income positions of United States banks over thirty years \citep{Begenau2015}. These works show that the real exposure of a bank depends on behavioural assumptions about deposits and on hedges, and that a standardised measure is a deliberate simplification. They do not treat a unit of account indexed to prices, because the markets they study have none, and so they do not tell what the separate treatment of such a unit adds to a standardised measure.

\subsection{Attribution of a difference to its components}
\label{sec:rw-attribution}

The measurement of this paper attributes the difference between two computations to the components of the treatment. Oaxaca decomposes a mean difference between two groups into a part explained by differences in characteristics and a part explained by differences in coefficients, in an analysis of wage differentials, and the decomposition depends on the order in which the parts are taken \citep{Oaxaca1973}. Shapley's value, defined for cooperative games, assigns to each player the average of its marginal contribution over all the orders in which the players can join, and it is the unique rule that is efficient, symmetric, additive and gives nothing to a dummy player \citep{Shapley1953}. Shorrocks uses it for the decomposition procedures of distributional analysis, because it gives a unified framework in which the contribution of a component does not depend on an arbitrary order \citep{Shorrocks2013}. Applied here the players are the five switches, the value of a coalition is a measure computed by the engine, and the efficiency property guarantees that the five contributions add up to the difference between the Annex 1 and the single-currency measures; the dummy property is a check, because a switch that cannot change a measure must receive zero, and Section~\ref{sec:theory-shapley} states two cases in which it must.

\subsection{Where the proposed measurement stands}
\label{sec:rw-stands}

The published methods cover the standardised measure for currencies in a table, the curves and models of indexed bonds, the exposure of real banks and the attribution of a difference to components. None combines them for an indexed unit of account in a standardised measure. The supervisor's text defines the treatment of the UF but does not decompose it, and the Basel standard cannot be applied to the peso or the UR without a national table. The proposed measurement takes the supervisor's text as the specification, implements it as five switches over an open engine, runs the Basel positive-part aggregation and the textbook gap and duration measures on the same flows as comparators, and attributes the difference between the separate and the single-currency treatments with the exact Shapley value. Its limits are stated in Section~\ref{sec:limitations}: it measures a rule, not an exposure; the bank is synthetic; and the engine has no second implementation of Annex 1 to reconcile with.
